\begin{definition}
Fix a finite loopless multigraph \(G\), retaining its exact edge type, a
palette of \(k\) colors, a distinguished edge \(e_0\), and a natural number
\(n\).  A \emph{technical induction context up to rung \(n\)} is a record
with the ten fields described below.

Every field quantifies independently over an arbitrary input partial
coloring \(\rho\), an arbitrary displayed ordered tree-sequence \(T\), an
arbitrary generating coloring \(\varphi\), and an
\[
  \mathcal E:\noderef{ExtendedTashkinovTree}(G,k,e_0,\rho,T,\varphi).
\]
Write \(s\) for the rung number of \(\mathcal E\), write
\(a_s=\operatorname{Fin.last}(s)\) for the last index in
\(\operatorname{Fin}(s+1)\), and, for \(i\in\operatorname{Fin}(s+1)\), put
\[
 L_i=\mathcal E.\mathrm{series.data.tree}(i.\mathrm{castSucc}),\qquad
 D_i=\noderef{TashkinovStoredColors}(\mathcal E.\mathrm{series.data},i),
 \qquad \varphi_i=\mathcal E.\mathrm{series.data.coloring}(i).
\]
Thus \(L_{a_s}\) is the exact terminal ladder tree.  Separately, put
\[
 F_s=\mathcal E.\mathrm{series.data.tree}
       (\operatorname{Fin.last}(s+1));
\]
this is the exact final series tree, whose index lies in
\(\operatorname{Fin}(s+2)\).  Also write \(\Theta_i,u_i,\gamma_i,\delta_i\)
for the extension kind, chosen vertex, gamma color, and delta color stored
at data index \(i\), and write \(S_i^{\rm raw}\) for
\(\mathcal E.\mathrm{series.data.stored}(i)\).  Membership in a displayed
tree always means membership in its
\(\noderef{TreeSequenceVertexSet}\).

The six lower-rung fields are universally quantified as above and have the
common hypotheses
\[
  0<s<n \quad\hbox{and}\quad \noderef{MaximumProperty}(\mathcal E).
\]
Their conclusions are the following.

\begin{enumerate}
\item[\texttt{lowerOne}.]
The displayed set \(V(T)\) is
\(\noderef{ColoredVertexSetElementary}\) with respect to the generating
coloring \(\varphi\).

\item[\texttt{lowerTwo}.]
For every coloring \(\sigma\), if \(\sigma\) is
\(\noderef{IsProperEdgeColoring}\), leaves exactly \(e_0\) uncolored in the
sense of \(\noderef{IsOneEdgeUncolored}\), and is
\(\noderef{StableColoring}\) on \(F_s\) relative to
\((D_{a_s},\varphi_{a_s})\), then for all colors \(\alpha,\beta\),
\[
 \alpha\in\noderef{TreeMissingColors}(\sigma,F_s)
 \quad\hbox{or}\quad
 \beta\in\noderef{TreeMissingColors}(\sigma,F_s)
 \quad\Longrightarrow\quad
 \noderef{TreeInterchangeable}(G,F_s,\sigma,\alpha,\beta).
\]
This is the direct universal interchangeability assertion on the exact
tree \(F_s\); in particular, the cited definition counts isolated
two-color path components.

\item[\texttt{lowerThree}.]
This field is a conjunction of two assertions.  First, for every
\(q\in\operatorname{Fin}(s)\), vertex \(v\), membership witness
\(h:v\in V(T)\), and \(j\in\operatorname{Fin}(s+1)\), if
\[
 \Theta_{q^+}=\mathrm{PE},\qquad u_{q^+}=v,\qquad
 \noderef{ETTFirstAppearanceIndex}(\mathcal E,v,h)=j+1,
\]
then every proper coloring \(\sigma\) leaving exactly \(e_0\) uncolored
which is stable on \(L_{q^+}\) relative to
\((D_{q^+},\varphi_{q^+})\) is
\(\noderef{StableColoringBeforeVertex}\) before \(v\) in the displayed
tree \(T\), using the same witness \(h\), relative to
\((D_j,\varphi_j)\).  In addition, under the explicit condition \(0<j\),
the same \(\sigma\) is stable on the predecessor ladder tree \(L_j\)
relative to \((D_j,\varphi_j)\).

Second, for all distinct \(q,r\in\operatorname{Fin}(s)\), all distinct
vertices \(v_q,v_r\), and membership witnesses
\(h_q:v_q\in V(T)\), \(h_r:v_r\in V(T)\), if
\[
 \Theta_{q^+}=\Theta_{r^+}=\mathrm{PE},\qquad
 u_{q^+}=v_q,\quad u_{r^+}=v_r,
\]
and the two dependent first-appearance indices computed with \(h_q,h_r\)
are equal, then
\[
  S_{q^+}^{\rm raw}\cap S_{r^+}^{\rm raw}=\varnothing.
\]

\item[\texttt{lowerFour}.]
For all \(v\in V(G)\) and colors \(\gamma,\delta\), if
\[
 \Theta_{a_s}=\mathrm{PE},\qquad u_{a_s}=v,\qquad
 \gamma_{a_s}=\gamma,\qquad\delta_{a_s}=\delta,
\]
then under every proper exactly-\(e_0\)-uncolored coloring \(\sigma\)
stable on \(L_{a_s}\) relative to
\((D_{a_s},\varphi_{a_s})\), every \(w\in V(L_{a_s})\) satisfies
\[
 \operatorname{inl}(w)\in
   \noderef{TwoColorComponent}(G,\sigma,\gamma,\delta,
      \operatorname{inl}(v))
 \quad\Longleftrightarrow\quad w=v.
\]

\item[\texttt{lowerFive}.]
Let \(\sigma\) be any proper exactly-\(e_0\)-uncolored coloring stable on
\(L_{a_s}\) relative to \((D_{a_s},\varphi_{a_s})\).  Let \(\delta\) be a
\(\noderef{DefectiveColor}\) of \(V(L_{a_s})\) under \(\sigma\), and let
\(v\) satisfy
\(\noderef{BoundaryColorVertex}(G,\sigma,V(L_{a_s}),\delta,v)\).
Assume there exists \(w\) satisfying the same exact boundary-color-vertex
predicate and
\(\noderef{TreeVertexOrder}(L_{a_s},w,v)\).  Then, for every certified
membership witness \(h:v\in V(T)\), every
\(q\in\operatorname{Fin}(s)\), and every vertex \(v_i\), if
\[
 \Theta_{q^+}=\mathrm{PE}\ \hbox{or}\
 \Theta_{q^+}=\mathrm{SE},\qquad
 u_{q^+}=v_i,\qquad
 \noderef{ETTFirstAppearanceIndex}(\mathcal E,v,h)\leq q+1,
\]
then
\[
 v=v_i\quad\hbox{or}\quad\noderef{TreeVertexOrder}(T,v,v_i).
\]

\item[\texttt{lowerSix}.]
Every proper exactly-\(e_0\)-uncolored coloring \(\sigma\) stable on
\(L_{a_s}\) relative to \((D_{a_s},\varphi_{a_s})\) satisfies
\(\noderef{ModColoring}(\mathcal E,\sigma)\).
\end{enumerate}

The four current-rung fields again quantify independently over arbitrary
\(\rho,T,\varphi,\mathcal E\), but their common hypotheses are now
\[
  s=n \quad\hbox{and}\quad \noderef{MaximumProperty}(\mathcal E).
\]
They are the following four certificates, stated at this exact rung.

\begin{enumerate}
\item[\texttt{currentThree}.]
This field is a conjunction.  For every
\(q\in\operatorname{Fin}(s)\), vertex \(v\), witness \(h:v\in V(T)\), and
\(j\in\operatorname{Fin}(s+1)\), the hypotheses
\[
 \Theta_{q^+}=\mathrm{PE},\qquad u_{q^+}=v,\qquad
 \noderef{ETTFirstAppearanceIndex}(\mathcal E,v,h)=j+1
\]
imply that every proper exactly-\(e_0\)-uncolored coloring \(\sigma\)
stable on \(L_{q^+}\) relative to
\((D_{q^+},\varphi_{q^+})\) is stable before \(v\) in \(T\), with witness
\(h\), relative to \((D_j,\varphi_j)\); and, only when \(0<j\), it is also
stable on \(L_j\) relative to \((D_j,\varphi_j)\).
Moreover, for distinct \(q,r\in\operatorname{Fin}(s)\), distinct vertices
\(v_q,v_r\), and witnesses \(h_q:v_q\in V(T)\),
\(h_r:v_r\in V(T)\), if both slots are PE, have respective chosen vertices
\(v_q,v_r\), and their dependent first-appearance indices are equal, then
\(S_{q^+}^{\rm raw}\cap S_{r^+}^{\rm raw}=\varnothing\).

\item[\texttt{currentFour}.]
For all \(v\in V(G)\) and colors \(\gamma,\delta\), if
\[
 \Theta_{a_s}=\mathrm{PE},\qquad u_{a_s}=v,\qquad
 \gamma_{a_s}=\gamma,\qquad\delta_{a_s}=\delta,
\]
then under every proper exactly-\(e_0\)-uncolored coloring \(\sigma\)
stable on \(L_{a_s}\) relative to
\((D_{a_s},\varphi_{a_s})\), every \(w\in V(L_{a_s})\) satisfies
\[
 \operatorname{inl}(w)\in
   \noderef{TwoColorComponent}(G,\sigma,\gamma,\delta,
      \operatorname{inl}(v))
 \quad\Longleftrightarrow\quad w=v.
\]

\item[\texttt{currentFive}.]
Let \(\sigma\) be any proper exactly-\(e_0\)-uncolored coloring stable on
\(L_{a_s}\) relative to \((D_{a_s},\varphi_{a_s})\).  Let \(\delta\) be a
\(\noderef{DefectiveColor}\) of \(V(L_{a_s})\), let \(v\) be a
\(\delta\)-boundary-color vertex of \(V(L_{a_s})\), and assume that there
exists another such boundary-color vertex \(w\) with
\(\noderef{TreeVertexOrder}(L_{a_s},w,v)\).  Then, for every certified
membership witness \(h:v\in V(T)\), every
\(q\in\operatorname{Fin}(s)\), and every vertex \(v_i\), if
\[
 \Theta_{q^+}=\mathrm{PE}\ \hbox{or}\
 \Theta_{q^+}=\mathrm{SE},\qquad
 u_{q^+}=v_i,\qquad
 \noderef{ETTFirstAppearanceIndex}(\mathcal E,v,h)\leq q+1,
\]
then \(v=v_i\) or \(\noderef{TreeVertexOrder}(T,v,v_i)\).

\item[\texttt{currentSix}.]
Every proper exactly-\(e_0\)-uncolored coloring \(\sigma\) stable on
\(L_{a_s}\) relative to \((D_{a_s},\varphi_{a_s})\) satisfies
\(\noderef{ModColoring}(\mathcal E,\sigma)\).
\end{enumerate}

Thus the lower fields package assertions (i)--(vi) strictly below the
bound, while the current fields package only assertions (iii)--(vi) at the
bound; assertions (i) and (ii) at rung \(n\) are deliberately not fields of
this context.
\end{definition}
